\chapter{Alternative Regularity via Non-Excessive Sweepouts}\label{ch:alternative-regularity}

In this chapter, we combine nested sweepouts with the non-excessive framework of Chodosh--Liokumovich--Spolaor~\cite{CLS22} and apply Wickramasekera's regularity theorem~\cite{Wic14} to obtain minimal hypersurfaces in closed Riemannian manifolds.

\section{Non-Excessive Sweepouts}\label{sec:non-excessive-prelim}
An alternative approach to min--max theory, developed by Chodosh--Liokumovich--Spolaor~\cite{CLS22}, is based on the notion of \emph{non-excessive sweepouts}. The non-excessive condition provides additional structural control that simplifies the regularity analysis.

We use the notation from Section~\ref{sec:sweepouts}. Recall that a sweepout $\{\Omega_t\}_{t\in[0,1]}$ of $M$ is a continuous family of Caccioppoli sets with $\Omega_0 = \emptyset$ and $\Omega_1 = M$ (Definition~\ref{def:sweepout}). The width is $W = W(M)$ as in Definition~\ref{def:width}.

\begin{definition}[Optimal nested volume parametrized sweepout]\label{def:onvp}
A sweepout $\{\Phi(x) = \partial \Omega(x)\}$ is called
\begin{itemize}
    \item \emph{optimal} if $\mathcal{F}(\Phi) = W$;
    \item \emph{nested} if $\Omega(x_1) \subset \Omega(x_2)$ for all $0\leq x_1 \leq x_2\leq 1$ (cf.\ Definition~\ref{def:nested-sweepout});
    \item \emph{volume parametrized} if $\Vol(\Omega(x)) = x$ for every $x\in [0,1]$.
\end{itemize}
We call such a sweepout an \emph{optimal nested volume parametrized (ONVP) sweepout}.
\end{definition}

\begin{definition}[Critical set]\label{def:critical-set}
Given a sweepout $\Phi$, define
\[
M(x) = \limsup_{r \to 0} \{\mathbf{M}(\Phi(y)): |y-x| < r\}.
\]
If $\Phi$ is an optimal sweepout, the \emph{critical domain of $\Phi$} is
\[
\mathfrak{m}(\Phi)=\{ x \in [0,1]: M(x) =W \}.
\]
We denote by $\mathfrak{m}_L(\Phi)$ and $\mathfrak{m}_R(\Phi)$ the sets of left and right critical points.
\end{definition}

\subsection{Excessive Intervals and Non-Excessive Points}

We introduce the concept of excessive intervals, which identify portions of a sweepout that can be replaced by lower-mass competitors.

\begin{definition}[Excessive intervals and points]\label{def:excessive-interval}
Suppose $\{\Phi(x) = \partial\Omega(x)\}$ is a sweepout. Given a connected interval $I \subset [0,1]$
(which may be open, closed, or half-open), we say that $\{\Phi^I(x) = \partial\Omega^I(x)\}_{x\in \bar{I}}$
is an \emph{$I$-replacement family for $\Phi$} if
$\Omega^I(a) = \Omega(a)$, $\Omega^I(b) = \Omega(b)$ at the endpoints, and for all $x \in I$,
\[
\limsup_{I \ni y \to x} \mathbf{M}(\Phi^I(y)) < W,
\]
where $W := \mathcal{F}(\Phi)$.

We say that a connected interval $I$ is an \emph{excessive interval for $\Phi$} if there exists an $I$-replacement family
for $\Phi$. A point $x$ is called \emph{left excessive for $\Phi$} (resp.\ \emph{right excessive}) if there exists an excessive interval $I$ with
$(x-\varepsilon, x] \subset I$ (resp.\ $[x, x+\varepsilon) \subset I$) for some $\varepsilon > 0$.
\end{definition}

\subsection{Non-Excessive Sweepouts and Stationary Varifolds}\label{subsec:stationary-varifolds-non-excessive}

The fundamental result of Chodosh--Liokumovich--Spolaor is the existence of sweepouts where critical points are not excessive:

\begin{theorem}[Existence of non-excessive min-max hypersurface~{\cite[Theorem~3.1]{CLS22}}]\label{thm:non-excessive-existence}
Let $(M^{n+1},g)$ be a closed Riemannian manifold with $2\leq n\leq 6$. There exists an (ONVP) sweepout $\Psi$ such that every $x\in \mathfrak{m}_L(\Psi)$ is not left excessive and every $x \in \mathfrak{m}_R(\Psi)$ is not right excessive.
\end{theorem}

A sweepout satisfying the conclusion of Theorem~\ref{thm:non-excessive-existence} is called a \emph{non-excessive sweepout}. We denote by $\Snm$ the collection of all such sweepouts.

\begin{proposition}[Stationary varifolds from optimal sweepouts]\label{thm:CLS-stationary}
Let $(M^{n+1},g)$ be a closed Riemannian manifold with $2\leq n\leq 6$, and let $\Phi$ be an optimal sweepout with $\mathcal{F}(\Phi) = W$ (the min-max width). Then there exists a stationary $n$-varifold $V$ in $M$ with $\mathbf{M}(V) \leq W$.
\end{proposition}

\begin{proof}
Each slice $\Phi(x) = \partial\Omega(x)$ induces a varifold $V_x = |\partial^*\Omega(x)|$. Let $\{x_i\}$ be any sequence with $\mass(\Phi(x_i)) \to W$. By varifold compactness, some subsequence satisfies $V_{x_i} \weakstar V$ with $\|V\|(M) \leq W$. We claim some such limit $V$ is stationary.

Suppose not: every varifold limit of sequences $\{V_{x_i}\}$ with $\mass(\Phi(x_i)) \to W$ is non-stationary. Consider $A_\eta = \{x : \mass(\Phi(x)) \geq W - \eta\}$, which is compact. By assumption, for each $x \in A_\eta$, any limit point of $\{V_y\}_{y \to x}$ is non-stationary. The tightening flow of \cite{CL03} strictly decreases mass near non-stationary varifolds, so by compactness of $A_\eta$, there exists $\varepsilon > 0$ (depending on $\eta$) such that tightening decreases mass by at least $\varepsilon$ on $A_\eta$.

Choose $\eta < \varepsilon$. After tightening, slices from $A_\eta$ have mass $\leq W - \varepsilon < W - \eta$, and slices outside $A_\eta$ already had mass $< W - \eta$. Thus the tightened sweepout $\tilde{\Phi}$ satisfies $\sup_x \mass(\tilde{\Phi}(x)) < W$, contradicting the definition of width.
\end{proof}

This is the standard pull-tight argument of~\cite{CL03}. Note that the non-excessive condition is not used here; it will be essential later for establishing stability and controlling mass cancellation.

\begin{remark}
The varifold $V$ obtained from Proposition~\ref{thm:CLS-stationary} is stationary, but its \emph{integrality} is not automatically guaranteed. This is the key issue that distinguishes the non-excessive framework from the classical Almgren--Pitts approach:
\begin{itemize}
    \item In Almgren--Pitts theory, one works with integral currents $\mathbf{Z}_n(M;\mathbb{Z}_2)$ from the outset, so integrality is built into the framework.
    \item In the non-excessive framework, one works with varifolds, and integrality must be established separately.
    \item Both approaches use pull-tight to obtain stationarity, but pull-tight alone does not produce integrality.
\end{itemize}
\end{remark}

\subsection{Homotopic Minimizers and Non-Minimizing Points}

To analyze the geometric properties of min-max varifolds, we introduce the notion of homotopic
minimizers and points where this property fails.

\begin{definition}[Homotopic minimizers]\label{def:homotopic-minimizer}
Let $\Sigma = \bdy\Omega$ be a Caccioppoli set in $M$, and let $U \subset M$ be an open set. We say
that $\Sigma$ is:
\begin{itemize}
    \item \emph{inner homotopic minimizing in $U$} if there is no Caccioppoli set $E \subset \Omega$
    with $E \triangle \Omega \subset U$ and $\per(E, U) < \per(\Omega, U)$ that can be connected to
    $\Omega$ by a continuous family of Caccioppoli sets in $U$;
    \item \emph{outer homotopic minimizing in $U$} if there is no Caccioppoli set $E \supset \Omega$
    with $E \triangle \Omega \subset U$ and $\per(E, U) < \per(\Omega, U)$ that can be connected to
    $\Omega$ by a continuous family of Caccioppoli sets in $U$;
    \item \emph{one-sided homotopic minimizing in $U$} if it is either inner or outer homotopic
    minimizing in $U$.
\end{itemize}
\end{definition}

\begin{definition}[Non-homotopic-minimizing points]\label{def:hnm}
For a varifold $V$ arising from a min-max procedure, we define the set of \emph{non-homotopic-minimizing
points} by
\[
    \mathfrak{h}_{\mathrm{nm}}(V) := \left\{P \in \spt\|V\| : \begin{array}{l}
        \text{for all $r > 0$ small, $\spt\|V\| \cap B_r(P)$ is not one-sided} \\
        \text{homotopic area minimizing on either side in $B_r(P)$}
    \end{array}\right\}.
\]
\end{definition}

Intuitively, $\mathfrak{h}_{\mathrm{nm}}(V)$ consists of points where the varifold $V$ cannot be
approximated by one-sided area minimizers in any neighborhood. These notions will be used in Section~\ref{sec:alpha-verification} to verify the $\alpha$-structural hypothesis required by Wickramasekera's regularity theorem.

\subsection{The Integrality Issue}\label{subsec:integrality-issue}

Following~\cite{CLS22}, we distinguish two cases for a critical point $x_0 \in \mathfrak{m}(\Phi)$ of a non-excessive sweepout $\Phi$:

\begin{definition}[Mass cancellation]\label{def:mass-cancellation}
Let $\Phi$ be a non-excessive (ONVP) sweepout with width $W$, and let $x_0 \in \mathfrak{m}(\Phi)$ be a critical point. We say that:
\begin{itemize}
    \item \emph{No mass cancellation} occurs at $x_0$ if $\mathbf{M}(\Phi(x_0)) = W$.
    \item \emph{Mass cancellation} occurs at $x_0$ if $\mathbf{M}(\Phi(x_0)) < W$.
\end{itemize}
\end{definition}

Without mass cancellation, the limiting varifold $V$ equals $|\partial\Omega(x_0)|$ with no mass loss. With mass cancellation, $V$ may have larger mass than $\mathbf{M}(\partial\Omega(x_0))$ due to sheets with opposite orientations canceling in the current while contributing to varifold mass.

In this chapter, we establish integrality and regularity without mass cancellation via Wickramasekera's theorem. The case with mass cancellation remains open.



\section{Main Result}

\begin{theorem}[Existence of minimal hypersurfaces via non-excessive min-max]\label{thm:non-excessive-main}
    Let $(M^{n+1}, g)$ be a closed Riemannian manifold with $2 \leq n \leq 6$. Then there exists a
    smooth closed embedded minimal hypersurface in $M$.
\end{theorem}

{\color{red}
\begin{conjecture}[Integrality with mass cancellation]\label{conj:mass-cancellation-integrality}
    If the tightened stationary varifold is integral even in the presence of mass cancellation, then
    Wickramasekera's regularity theorem remains applicable, and the conclusion of
    Theorem~\ref{thm:non-excessive-main} extends to this case as well.
\end{conjecture}
}


\section{Integrality without Mass Cancellation}

\begin{theorem}[Integrality without mass cancellation]\label{thm:integrality-mult-one}
    Let $(M^{n+1},g)$ be a closed Riemannian manifold with $n \geq 2$, and let $V$ be a
    stationary varifold arising from a non-excessive sweepout via the construction of
    Theorem~\ref{thm:CLS-stationary}. If $\Theta^n(V,x)=1$ for $\mathbf{M}(V)$-almost every
    $x\in\spt(V)$, then $V$ is integral and rectifiable.
\end{theorem}

\subsection{Case 1: No Mass Loss}

The reduced boundary $\bdy\Omega$ of a Caccioppoli set $\Omega$ induces a varifold
$V = \mathcal{H}^n \llcorner \bdy\Omega \otimes \delta_{T_x\bdy\Omega}$.

The key tool is the following proposition from the appendix of De~Lellis--Tasnady~\cite{DLT13}:

\begin{proposition}[Varifolds vs Caccioppoli set limits~{\cite[Proposition A.1]{DLT13}}]\label{prop:DLT-integrality}
    Let $\{\Omega^k\}$ be a sequence of Caccioppoli sets and $U$ an open subset of $M$. Assume that
    \begin{itemize}
        \item[(i)] $D\mathbb{1}_{\Omega^k} \to D\mathbb{1}_\Omega$ in the sense of measures in $U$;
        \item[(ii)] $\per(\Omega^k, U) \to \per(\Omega, U)$
    \end{itemize}
    for some Caccioppoli set $\Omega$. Denote by $V^k$ and $V$ the varifolds induced by $\bdy\Omega^k$
    and $\bdy\Omega$ respectively, where $\bdy\Omega$ denotes the reduced boundary. Then
    $V^k \to V$ in the sense of varifolds.
\end{proposition}

\begin{proof}[Proof of Theorem~\ref{thm:integrality-mult-one}, Case 1]
By Theorem~\ref{thm:CLS-stationary}, the varifold $V$ arises as the weak limit of varifolds $V_i$
associated to slices $T_i = \phi(t_i)$ with $\mathbf{M}(T_i) \to \mathcal{F}(\phi)$.

We identify each slice $T_i \in \mathbf{Z}_n(M;\mathbb{Z}_2)$ with its corresponding Caccioppoli set
$\Omega_i$ (see Proposition~\ref{prop:cacc-cycle}). The
varifold $V_i$ is then the varifold induced by the reduced boundary $\bdy\Omega_i$, given by
\[
    V_i = \mathcal{H}^n \llcorner \bdy\Omega_i \otimes \delta_{T_x\bdy\Omega_i}.
\]

\textbf{Step 1: BV convergence.} By the compactness theorem for BV functions (Caccioppoli sets), the
sequence $\{\Omega_i\}$ has a subsequence (still denoted $\{\Omega_i\}$) such that
$D\mathbb{1}_{\Omega_i} \to D\mathbb{1}_\Omega$ weakly in the sense of measures, for some
Caccioppoli set $\Omega$.

\textbf{Step 2: Perimeter convergence.} We establish perimeter convergence using the no mass cancellation hypothesis.

By passing to a subsequence, we may assume $t_i \to t^*$ for some $t^* \in [0,1]$. Since $\mathbf{M}(\phi(t_i)) \to W = \mathcal{F}(\phi)$ and the function $M(x) := \limsup_{r \to 0} \{\mathbf{M}(\phi(y)) : |y - x| < r\}$ is upper semicontinuous, we have $M(t^*) \geq W$. But $M(x) \leq W$ for all $x$ by definition of width, so $M(t^*) = W$, i.e., $t^* \in \mathfrak{m}(\phi)$.

By the \emph{no mass cancellation} hypothesis (Definition~\ref{def:mass-cancellation}), we have
\[
    \mathbf{M}(\phi(t^*)) = W.
\]
Since the sweepout $\phi$ is continuous in the flat norm topology (see Proposition~\ref{prop:sweepout-cycle}), and flat convergence implies $L^1$ convergence of characteristic functions, we have $\mathbb{1}_{\Omega_i} \to \mathbb{1}_{\Omega(t^*)}$ in $L^1(M)$. By uniqueness of BV limits, the Caccioppoli set $\Omega$ from Step~1 equals $\Omega(t^*)$.

Now we can establish perimeter convergence directly:
\begin{itemize}
    \item \textbf{Upper bound:} By the identification of currents and Caccioppoli sets,
    \[
        \per(\Omega_i, M) = \mathbf{M}(\phi(t_i)) \to W.
    \]
    \item \textbf{Lower bound:} By lower semicontinuity of perimeter under BV convergence,
    \[
        \per(\Omega, M) \leq \liminf_{i \to \infty} \per(\Omega_i, M) = W.
    \]
    \item \textbf{Equality:} Since $\Omega = \Omega(t^*)$ and $\mathbf{M}(\phi(t^*)) = W$ by no mass cancellation,
    \[
        \per(\Omega, M) = \per(\Omega(t^*), M) = \mathbf{M}(\phi(t^*)) = W.
    \]
\end{itemize}
Combining these, we obtain perimeter convergence:
\[
    \per(\Omega_i, M) \to \per(\Omega, M) = W.
\]

\textbf{Step 3: Application of Proposition~\ref{prop:DLT-integrality}.} By
Proposition~\ref{prop:DLT-integrality} (De~Lellis--Tasnady, Proposition A.1), since we have both
BV convergence and perimeter convergence, the varifold limit $V$ equals the varifold induced by
$\bdy\Omega$:
\[
    V = \mathcal{H}^n \llcorner \bdy\Omega \otimes \delta_{T_x\bdy\Omega}.
\]
This is precisely the varifold induced by an integral current (the Caccioppoli set $\Omega$), hence
$V$ is integral and rectifiable.
\end{proof}

\subsection{Case 2: Mass Cancellation}

In this case, we assume that $\mathbf{M}(\Phi(x_0)) < W$, so that mass cancellation occurs at the critical point $x_0$. As noted in Conjecture~\ref{conj:mass-cancellation-integrality}, establishing integrality in this case remains an open problem.

\begin{proposition}[Regularity assuming Conjecture~\ref{conj:mass-cancellation-integrality}]\label{prop:regularity-with-cancellation}
    Suppose Conjecture~\ref{conj:mass-cancellation-integrality} holds, so that the stationary varifold $V$ arising from a non-excessive sweepout is integral even in the presence of mass cancellation. Then the same regularity argument as in Case 1 applies, and $V$ is a smooth embedded minimal hypersurface.
\end{proposition}

\begin{remark}
The fundamental obstruction to proving Conjecture~\ref{conj:mass-cancellation-integrality} is that when $\mathbf{M}(V) < W$, we cannot apply the index bound $\mathrm{Index}(V) \leq 1$ from min-max, which prevents us from using Wickramasekera's regularity theory to establish integrality. This represents a genuine obstruction rather than merely a technical gap.
\end{remark}

\begin{proof}
If $V$ is integral, then by Proposition~\ref{prop:DLT-integrality}, there exists a Caccioppoli set $\Omega$ such that
\[
    V = \mathcal{H}^n \llcorner \partial^*\Omega \otimes \delta_{T_x\partial^*\Omega}.
\]
The reduced boundary $\partial^*\Omega$ is $(n-1)$-rectifiable by De Giorgi's theorem, hence $V$ is rectifiable and integral.

The regularity then follows from Wickramasekera's theorem (Theorem~\ref{thm:wickramasekera}) once we verify the $\alpha$-structural hypothesis, which can be done using the same deformation arguments as in Section~\ref{sec:alpha-verification} for the case without mass cancellation.
\end{proof}

\section{Regularity via Wickramasekera's Theorem}

We recall the main regularity result from Wickramasekera~\cite{Wic14}:

\begin{theorem}[Wickramasekera~{\cite{Wic14}}]\label{thm:wickramasekera}
    Let $V$ be a stable integral $n$-varifold in a closed $(n+1)$-dimensional Riemannian manifold
    with $2\leq n\leq 6$. Suppose that for some $\alpha \in (0, 1/2)$, $V$ satisfies the $\alpha$-structural hypothesis (Definition~\ref{defn:alpha-structural}).
    Then $\sing V = \emptyset$, i.e., $\spt(V)$ is a smooth embedded minimal hypersurface.
\end{theorem}

\begin{definition}[$\alpha$-structural hypothesis]\label{defn:alpha-structural}
For $\alpha \in (0, 1/2)$, a stationary integral varifold $V$ satisfies the \emph{$\alpha$-structural hypothesis} if \textbf{no singular point}
of $V$ has a neighborhood in which $V$ corresponds to a union of embedded $C^{1,\alpha}$ hypersurfaces-with-boundary
meeting \emph{only} along their common $C^{1,\alpha}$ boundary, and with multiplicity a constant positive integer
on each of the constituent hypersurfaces-with-boundary.
\end{definition}

\begin{remark}
By Wickramasekera's theorem, the following are equivalent for a stable integral $n$-varifold $V$:
\begin{itemize}
    \item[(a)] $V$ satisfies the $\alpha$-structural hypothesis for some $\alpha \in (0, 1/2)$.
    \item[(b)] $\sing V = \emptyset$ if $2 \leq n \leq 6$.
    \item[(c)] $\mathcal{H}^{n-1}(\sing V) = 0$ (vanishing $(n-1)$-dimensional Hausdorff measure of the singular set).
\end{itemize}
Thus, to prove regularity via Wickramasekera's theorem, it suffices to verify the $\alpha$-structural hypothesis for some $\alpha \in (0, 1/2)$.
\end{remark}

\subsection{Verification of the $\alpha$-Structural Hypothesis}\label{sec:alpha-verification}

\begin{theorem}[Verification of $\alpha$-structural hypothesis]\label{thm:alpha-hypo}
    Let $V$ be the stationary integral varifold from Theorem~\ref{thm:integrality-mult-one} arising from a
    non-excessive sweepout without mass cancellation. Then $V$ satisfies the $\alpha$-structural hypothesis.
\end{theorem}

\subsubsection{Case 1: No Mass Cancellation}

Assume the \emph{no cancellation} condition:
\begin{equation}
    \mass(\Phi(x_0)) = W,
\end{equation}
where $W$ is the min-max width and $x_0$ is a critical point.

In this case, the stationary integral varifold $V$ from Theorem~\ref{thm:integrality-mult-one} can be written as
$V = |\Sigma| = |\bdy\Omega|$ for some Caccioppoli set $\Omega$.

Our proof of Theorem~\ref{thm:alpha-hypo} under the no cancellation condition is similar to~\cite[Lemma 3.4]{CLS22}.

\begin{remark}[Why the case analysis is complete]\label{rem:case-completeness}
We explain why Cases 1a and 1b below are exhaustive, and why stability holds.

\textbf{(1) Completeness of cases.} Since $V = |\partial\Omega|$ in the no mass cancellation case, tangent cones at singular points are constrained by the Caccioppoli set structure. For stationary boundaries of Caccioppoli sets in $\mathbb{R}^{n+1}$ ($n \leq 6$), the only possible singular cones are:
\begin{itemize}
    \item \textbf{Case 1a:} $C = |H_1| + |H_2|$ (two half-hyperplanes meeting at angle $< \pi$) --- a ``wedge''
    \item \textbf{Case 1b:} $C = 2|H|$ (single half-hyperplane with multiplicity 2) --- a ``fold''
\end{itemize}
Higher multiplicities ($q \geq 3$) or more complex configurations (three or more sheets meeting) would require mass cancellation, which is impossible for Caccioppoli boundaries.

\textbf{(2) Stability.} The min-max construction gives $\mathrm{Index}(V) \leq 1$. Combined with the deformation arguments below showing at most one unstable direction, we conclude $\mathrm{Index}(V) = 0$ (stability) or $\mathrm{Index}(V) = 1$ (one-sided minimization), both of which lead to regularity via Wickramasekera~\cite{Wic14} or Liu~\cite{Liu19}.
\end{remark}

\begin{proof}[Proof of Theorem~\ref{thm:alpha-hypo}, Case 1]
Suppose the $\alpha$-structural hypothesis fails at $Z \in \sing{V}$. Then there exists $\rho > 0$ such that
\begin{equation}
    B_\rho(Z)\backslash\spt{\norm{V}} = \left(\bigsqcup_{i=1}^NE_i\right)\sqcup\left(\bigsqcup_{j=1}^NF_j\right),
\end{equation}
and we denote $\mathcal{E}:= \bigsqcup_{i=1}^NE_i$ (resp.\ $\mathcal{F}:= \bigsqcup_{j=1}^NF_j$) to represent those
Caccioppoli sets within $\overline{\Omega}$ (resp.\ $\overline{\Omega^c}$).

Without loss of generality, we only need to show that the $E_i$ are not homotopic minimizers. For each $i = 1, \dots, N$,
we denote by $C$ the blow-up cone of $E_i$ at $Z$. Then there are only two possibilities for the structure of the cone:

\begin{itemize}
    \item \textbf{Case 1a:} $C = |H_1|+|H_2|$ where both $H_1$ and $H_2$ are half-hyperplanes gluing along their
    common boundaries with the angle between them less than $\pi$.

    We denote $C = |\bdy\Omega_C|$ for some Caccioppoli set $\Omega_C$ of $\mathbb{R}^{n + 1}$. In this case, as
    in~\cite[Lemma 3.4]{CLS22}, we can choose $C^{1, \alpha}$-coordinates on $M$ around $Z$ such that $E_i = \Omega_C$
    in $B_\varepsilon(Z)$ so that $g_{ij}(Z) = \delta_{ij}$, which we can do since $g \in C^2$ and $\Sigma = |\bdy E_i|$
    is a $C^{1, \alpha}$-deformation of $C$ near $Z$ by assumption.

    Since the angle between $H_1$ and $H_2$ is less than $\pi$, by the optimal regularity theorem for one-sided
    minimizers from~\cite{Liu19}, we can find a Caccioppoli set $E_C \subset \Omega_C$ so that
    $E_C\triangle \Omega_C\subset B_1 \subset \mathbb{R}^{n + 1}$ and
    \begin{equation}
        \per_{\mathbb{R}^{n + 1}}(E_C, B_1) \leq \per_{\mathbb{R}^{n + 1}}(\Omega_C, B_1) - \delta
    \end{equation}
    for some $\delta > 0$. We set
    \begin{equation}
        E(x) := \begin{cases}
            (E_i\backslash B_\varepsilon(Z))\cup(|x|E_C\cap B_\varepsilon(Z)) & x < 0\\
            E_i & x = 0.
        \end{cases}
    \end{equation}
    We have that
    \begin{equation}
        \per_g(E(x)) - \per_g(\Omega) = -|x|^n\delta(1 + o(1))
    \end{equation}
    as $x \to 0$. This shows that $\Sigma$ is not inner homotopic minimizing.

    \item \textbf{Case 1b:} $C = 2|H|$ where $H$ is a half-hyperplane. In this case, we define
    \begin{equation}
        E(x) := \begin{cases}
            E_i\backslash B_{|x|}(Z) & x < 0\\
            E_i & x= 0.
        \end{cases}
    \end{equation}
    Apparently, for any $r > 0$,
    \begin{equation}
        \mass_{B_1}(H\backslash B_r(0)) = \mass_{B_1}(H) - \frac{1}{2}\omega_n r^n.
    \end{equation}
    Moreover, since $E_i$ is also a two-valued $C^{1, \alpha}$ graph over $H$ and $|H|$ is the blow-up cone of
    $|\bdy E_i|$ at $Z$, we know that $\per_g(E_i, B_{|x|}(Z)) = o(1)|x|^n$. Thus, we can again conclude that
    \begin{equation}
        \per_g(E(x)) - \per_g(\Omega) = -|x|^n\delta(1 + o(1)),
    \end{equation}
    where $\delta = \omega_n/2$.
\end{itemize}

We can define the corresponding $F(x)$ for $F_j$ by changing $x < 0$ to $x > 0$. Thus, we can know that
$Z \in \mathfrak{h}_{nm}(\Sigma)$, where $\mathfrak{h}_{nm}$ denotes the set of points that are not homotopic minimizing
in the sense of~\cite{CLS22}:
\begin{equation}
    \mathfrak{h}_{nm}(V) :=
  \left\{P \in \spt{\norm{V}}: \begin{array}{l}
      \text{for all $r > 0$ small, $\spt\norm{V}\cap B_r(P)$ is not one-sided}\\
      \text{homotopic area minimizing on either side in $B_r(P)$}
  \end{array}\right\}.
\end{equation}

Since $Z \in \mathfrak{h}_{nm}(V)$, by~\cite[Proposition 3.1]{CLS22}, we know that there is $\rho$ sufficiently small
such that for every $X \notin B_\rho(Z)$, there exists $r > 0$ such that $\spt{\norm{V}}$ is one-sided minimizing in
$B_r(X)$. However, since the $\alpha$-structural hypothesis fails, we can take $X \in \sing{V}$ with $X \notin B_\rho(Z)$
such that $\spt{\norm{V}}\cap B_r(X)$ has $(n - 1)$-dimensional singularity, which contradicts the optimal regularity
for one-sided minimizers from~\cite{Liu19}.

\begin{remark}[Dimension of the singular set]\label{rem:sing-dimension}
The proof above uses the existence of singular points $X \in \sing V$ outside $B_\rho(Z)$. We explain why such points exist.

If the $\alpha$-structural hypothesis fails at $Z$, then by Cases 1a and 1b, the tangent cone at $Z$ has an $(n-1)$-dimensional ``spine'' $L$---the common boundary of the half-hyperplanes. Since $V$ approximates this cone in the Hausdorff sense near $Z$, the singular set $\sing V$ must contain an $(n-1)$-dimensional subset near $Z$, not just an isolated point. In particular, for any $\rho > 0$, there exist singular points $X \in \sing V \setminus B_\rho(Z)$ where the same structural analysis applies, leading to the contradiction in the proof.
\end{remark}

Therefore, the $\alpha$-structural hypothesis must hold.
\end{proof}

The min-max construction gives $\text{index}(V) \leq 1$. This index bound holds \emph{a priori} before regularity is established, avoiding any circular reasoning.

Thus if $V$ is stable (i.e., $\mathrm{index}(V) = 0$), then
Wickramasekera's regularity theorem~\cite{Wic14} guarantees the optimal regularity for $V$. And in the case when
$\text{index}(V) = 1$, $V$ is one-sided area minimizing, thus we can again use the optimal regularity for one-sided
minimizers from~\cite{Liu19}.

\subsubsection{Case 2: Mass Cancellation}

Assume the \emph{cancellation} condition: $\mass(\Phi(x_0)) < W$ throughout this section.

In this case, strict mass cancellation occurs: there exists a cancellation point $q \in \text{reg}\, V$ such that
\[
    \per(\Omega, B_\varepsilon(q)) < |V|(B_\varepsilon(q))
\]
for all sufficiently small $\varepsilon > 0$, where $\Omega = \Omega(x_0)$ and $V$ can be written as
$V = \sum_i\kappa_i|\Sigma_i|$ for integer multiplicities $\kappa_i$.

\begin{remark}[Integrality assumption]
The representation $V = \sum_i\kappa_i|\Sigma_i|$ with \emph{integer} multiplicities assumes integrality, which is precisely what Conjecture~\ref{conj:mass-cancellation-integrality} seeks to establish. The analysis below should be understood as conditional on this conjecture. The fundamental obstruction to proving integrality in the mass cancellation case is that without $\mathbf{M}(V) = W$, we cannot apply the index bound to deduce regularity.
\end{remark}

The key difference from Case 1 is that $\Sigma$ need not be globally representable as $\bdy\Omega$; instead,
it may consist of multiple components with different multiplicities. Nevertheless, the cancellation condition
allows us to control deformations away from the cancellation point $q$.

\begin{definition}[Competitors in the cancellation case]\label{def:competitor}
Given $\varepsilon, \tau > 0$ with $\Sigma = \spt(V) \subset \overline{\Omega}$, we call a Caccioppoli set
$\Omega'$ a \emph{$(q, \varepsilon, \tau, \Sigma, \Omega)$-competitor} if
\[
    (\Omega \setminus B_\tau(\Sigma)) \cup (B_\varepsilon(q) \setminus \Sigma) \subset \Omega' \subsetneq \Omega \setminus \Sigma.
\]
A $(q, \varepsilon, \tau, \Sigma, \Omega)$-competitor $\Omega'$ is called a \emph{minimizing competitor}
if its perimeter is strictly less than the perimeter of any $(q, \varepsilon, \tau, \Sigma, \Omega)$-competitor
$\Omega''$ with $\Omega' \subset \Omega''$.
\end{definition}

Intuitively, a competitor represents a Caccioppoli set that:
\begin{itemize}
    \item Agrees with $\Omega$ outside a $\tau$-neighborhood of $\Sigma$;
    \item Agrees with $\Omega$ near the cancellation point $q$ (in $B_\varepsilon(q)$);
    \item Has strictly smaller perimeter than $\Omega$ in the intermediate region.
\end{itemize}

The key technical result is that such competitors cannot exist when $x_0$ is not excessive:

\begin{proposition}[Non-existence of minimizing competitors~{\cite[Proposition~6.3]{CLS22}}]\label{prop:no-competitor}
Suppose $V$ arises from a non-excessive sweepout $\Phi$ with $x_i \nearrow x_0 \in \mathfrak{m}_L(\Phi)$
and $|\Phi(x_i)| \to V$, and suppose the cancellation condition $\mass(\Phi(x_0)) < W$ holds. Then for
every $\varepsilon > 0$, there exists $\tau > 0$ such that no minimizing
$(q, \varepsilon, \tau, \Sigma, \Omega)$-competitor exists.
\end{proposition}

\begin{proof}[Proof sketch]
Suppose a minimizing competitor $U$ exists. By the cancellation condition, we can find Caccioppoli sets
$\Omega'$ with $\per(\Omega') \leq W + \eta - \delta(\varepsilon)$ for small $\eta$ and $\delta(\varepsilon) > 0$
from the mass loss at $q$. The minimizing property of $U$ implies $\per(U) < W$.

For $\tau$ sufficiently small, a continuity argument (as in~\cite[Lemma 3.3]{CLS22}) provides a nested
family $\{E(x) : x \in [0,1]\}$ with $E(0) = U$, $E(1) = \Omega(x_0)$, and $\per(E(x)) < W$. Combining
this with the original sweepout produces a replacement family showing that $x_0$ is left excessive,
contradicting Theorem~\ref{thm:non-excessive-existence}.
\end{proof}

With the non-existence of minimizing competitors established, we can prove regularity (conditional on Conjecture~\ref{conj:mass-cancellation-integrality}):

\begin{proof}[Proof of Theorem~\ref{thm:alpha-hypo}, Case 2 (assuming integrality)]
By Proposition~\ref{prop:no-competitor}, there are no minimizing competitors for any $\varepsilon > 0$ and
suitable $\tau > 0$. This implies:

\textbf{Stability of $V$:} If any component $\Sigma_i$ had unstable regular part, one could perform a one-sided
minimization in $B_\tau(\Sigma_i) \setminus B_\varepsilon(q)$ to find a minimizing competitor,
contradicting Proposition~\ref{prop:no-competitor}. Thus each $\Sigma_i$ has stable regular part.

\textbf{$\mathfrak{h}_{\mathrm{nm}}(V) = \emptyset$:} By~\cite[Proposition~3.7]{CLS22}, the non-existence
of minimizing competitors implies that every point of $V$ is one-sided area minimizing in a small
neighborhood. Therefore, $\mathfrak{h}_{\mathrm{nm}}(V) = \emptyset$.

\textbf{Verification of $\alpha$-structural hypothesis:} If the $\alpha$-structural hypothesis fails at
some point $Z \in \sing{V}$, we can construct a local area-reducing deformation as in Case 1. However,
such a deformation would produce a minimizing competitor (by one-sided minimization), contradicting
Proposition~\ref{prop:no-competitor}.

Therefore, $V$ satisfies all hypotheses of Wickramasekera's theorem~\cite{Wic14}, implying smoothness.
\end{proof}

\subsection{Proof of the Main Theorem}

We are now ready to prove Theorem~\ref{thm:non-excessive-main}.

\begin{proof}[Proof of Theorem~\ref{thm:non-excessive-main}]
Let $(M^{n+1}, g)$ be a closed Riemannian manifold with $2 \leq n \leq 6$. The nested sweepout
framework constructed in Chapters~\ref{ch:morse-foliation}--\ref{ch:nested-sweepout} provides a continuous
family of hypersurfaces sweeping out $M$. Applying the non-excessive tightening procedure of
Chodosh--Liokumovich--Spolaor to this nested sweepout yields a non-excessive sweepout
$\Phi \in \mathcal{S}_{\text{nm}}$.

By Theorem~\ref{thm:CLS-stationary}, there exists a stationary varifold $V$ with $\mass(V) \leq \mathcal{F}(\Phi)$.

Assume there is no mass cancellation (i.e., $\Theta^n(V,x)=1$ for $\mass(V)$-almost every $x \in \spt(V)$).
By Theorem~\ref{thm:integrality-mult-one}, $V$ is integral and rectifiable.

By Theorem~\ref{thm:alpha-hypo}, $V$ satisfies the $\alpha$-structural hypothesis.

The non-excessive condition ensures that $V$ is stable. Therefore, all hypotheses of Wickramasekera's regularity
theorem (Theorem~\ref{thm:wickramasekera}) are satisfied. We conclude that $\sing V = \emptyset$, i.e., $\spt(V)$
is a smooth embedded minimal hypersurface.
\end{proof}

